% from ega4-2-fr.tex
\oldpage[IV-2]{170}
(ii) Comme $A\subset A_\lambda\subset R$ pour tout $\lambda$, $R$ est aussi
l'anneau total des fractions de chaque $A_\lambda$. Vu la caractérisation
des anneaux normaux ayant $R$ pour anneau total des fractions, donnée dans
(i), l'assertion (ii) résulte de Bourbaki, Alg. comm., chap. V, § 1, n$^o$
3, prop. 12.
\end{proof}

Ce lemme étant démontré, la démonstration de (6.14.1) procède en plusieurs
étapes.
I) \emph{Réduction au cas où $Y=\operatorname{Spec}(A)$,
$Y'=\operatorname{Spec}(A')$, $X=\operatorname{Spec}(B)$, $A$, $A'$ étant
des anneaux locaux noethériens, $A$ intègre, $B$ la clôture intégrale de
$A$.} --- Il s'agit de prouver que pour tout $x'\in X'$, l'anneau local
$\mathcal O_{X',x'}$ est intègre et intégralement clos; soient $x$, $y$,
$y'$ les images canoniques de $x'$ dans $X$, $Y$, $Y'$ respectivement; si
l'on pose $A=\mathcal O_{Y,y}$, $A'=\mathcal O_{Y',y'}$,
$B=\mathcal O_{X,x}$, les anneaux $\mathcal O_{X',x'}$ pour $x$, $y$, $y'$
fixés, sont les anneaux locaux aux idéaux premiers de l'anneau
$B'=B\otimes_A A'$ (I, 3.6.5), et il s'agit donc de prouver que $B'$ est un
anneau normal. Notons d'autre part que
$\operatorname{Spec}(A')\to\operatorname{Spec}(A)$ est un morphisme normal
par définition (6.8.1 et I, 3.6.5); on est donc ramené au cas où $Y$ et $Y'$
sont des schémas locaux, $X$ le spectre d'un anneau local intègre et
intégralement clos $B$. Désignons par $(B_\alpha)$ la famille des clôtures
intégrales des sous-$A$-algèbres de type fini de $B$; il est clair que $B$
est réunion de la famille filtrante croissante des $B_\alpha$. Comme le
foncteur $\varinjlim$ commute au produit tensoriel, $B'$ est donc isomorphe
à $\varinjlim B'_\alpha$, où l'on a posé
$B'_\alpha=B_\alpha\otimes_A A'$. Pour prouver que $B'$ est normal, il
suffira, en vertu de (5.13.6), de montrer que les anneaux $B'_\alpha$ sont
normaux et que, pour $\alpha\leqslant\beta$, toute composante irréductible
de $\operatorname{Spec}(B'_\beta)$ domine une composante irréductible de
$\operatorname{Spec}(B'_\alpha)$. Mais cette dernière propriété résulte de
l'hypothèse que $A'$ est un $A$-module plat et de (2.3.7, (ii)), puisque
$B_\alpha$ et $B_\beta$ sont intègres et que $B_\alpha\subset B_\beta$.

On peut donc se borner à démontrer que $B'=B\otimes_A A'$ est normal lorsque
$B$ est la clôture intégrale d'une $A$-algèbre intègre de type fini $C$;
dans ce cas, si l'on pose $C'=C\otimes_A A'$, le morphisme
$\operatorname{Spec}(C')\to\operatorname{Spec}(C)$ est normal (6.8.2);
comme $B'=B\otimes_C C'$, on voit qu'on peut remplacer $A$ et $A'$ par $C$
et $C'$ respectivement, donc supposer que $A$ est intègre et que $B$ est la
clôture intégrale de $A$. Enfin, le procédé du début permet de se borner au
cas où $A$ est local (tenant compte de ce que, si $B$ est la clôture
intégrale de $A$, $B_\mathfrak p$ est la clôture intégrale de
$A_\mathfrak p$ pour tout idéal premier $\mathfrak p$ de $A$).

